\documentclass[11pt,a4paper]{article}

\usepackage[utf8]{inputenc}
\usepackage[T1]{fontenc}
\usepackage{lmodern}
\usepackage{microtype}
\usepackage[margin=2.5cm]{geometry}
\usepackage{amsmath,amssymb,amsthm}
\usepackage{mathtools}
\usepackage{booktabs}
\usepackage{longtable}
\usepackage{array}
\usepackage{enumitem}
\usepackage{hyperref}
\usepackage{xcolor}
\usepackage{tcolorbox}
\usepackage{fancyhdr}
\usepackage{titlesec}
\usepackage{csquotes}

\hypersetup{
  colorlinks=true,
  linkcolor=blue!50!black,
  citecolor=blue!50!black,
  urlcolor=blue!50!black,
}

\pagestyle{fancy}
\fancyhf{}
\fancyhead[L]{\small Evidence Ledger for the V$_{600}$ Closure Programme}
\fancyhead[R]{\small Pre-peer-review preprint}
\fancyfoot[C]{\thepage}

\newtheorem{definition}{Definition}
\newtheorem{principle}{Principle}
\newtheorem{criterion}{Criterion}

\newcommand{\II}{\mathcal{I}}
\newcommand{\IIlip}{\mathcal{I}_{\mathrm{Lip}}}
\newcommand{\IImax}{\mathcal{I}_{\max}}
\newcommand{\Vsix}{V_{600}}
\newcommand{\Cphi}{C_\varphi}
\newcommand{\NH}{N_H}
\newcommand{\Q}{\mathbb{Q}}
\newcommand{\Z}{\mathbb{Z}}
\newcommand{\R}{\mathbb{R}}
\newcommand{\C}{\mathbb{C}}
\newcommand{\N}{\mathbb{N}}
\newcommand{\OK}{\mathcal{O}_K}

\newtcolorbox{scopebox}[1][]{
  colback=blue!2,
  colframe=blue!40!black,
  fonttitle=\bfseries,
  title=Scope,
  #1
}

\newtcolorbox{burdenbox}[1][]{
  colback=red!2,
  colframe=red!40!black,
  fonttitle=\bfseries,
  title=The burden,
  #1
}

\title{\textbf{From Shared Vocabulary to Independent Constraint}\\[0.4em]
\large An Evidence Ledger for the $V_{600}$ Closure Programme}

\author{Lee Smart\\
\small Institute of Vibrational Field Dynamics\\
\small \texttt{contact@vibrationalfielddynamics.org}}

\date{2026-05-28 \quad (pre-peer-review preprint, v1.0.0-rc1)}

\begin{document}
\maketitle

\begin{abstract}
\noindent
The $V_{600}$ closure programme assembles a single mathematical
substrate --- the icosian triad $(\II,\mathcal{G},\Cphi)$ on the
$600$-cell vertex set $\Vsix$ --- and applies it across mathematics,
cosmology, biology and number theory. A natural and serious critique
is that this recurrence is methodological rather than evidential:
once a vocabulary becomes the working language of a programme, it
will appear everywhere the programme looks. We agree this critique
is correct unless it can be shown, claim by claim, that the
substrate does work that an unconstrained free-parameter model would
not do.

This paper is the audit. We separate the programme's outputs into
four classes: exact mathematical certificates (Class A), constrained
physical derivations (Class B), data-facing empirical analyses
(Class C), and interpretive ontology (Class D). For each individual
claim we record the starting constraints, the role $\Vsix$ or
closure plays, the number of free parameters fitted versus inputs
fixed, the output, the comparison target, the reproducibility
artefact, and the falsification condition. The result is an
\emph{evidence matrix} (Appendix A), a \emph{claim-status table}
(Appendix B), and a \emph{falsification ledger} (Appendix C).

We do not claim that the result vindicates the programme. We claim
that it allows the programme's evidential standing to be assessed at
the granularity of individual claims rather than as a single
holistic gesture. Some claims (Class A) survive any methodological
critique because they are theorem-grade. Some (Class C) survive
because the substrate fixed parameters that were free in the
comparison model. Some (Class B) are conditional on named bridge
hypotheses that we list as open. Some (Class D) are interpretation
and are flagged as such.

The reader who concludes after this audit that the programme is
``mostly Class D dressed in Class A clothing'' will, we hope, be
able to point to the specific row of the ledger that supports that
verdict. The reader who concludes the opposite will be able to do
the same.
\end{abstract}

\bigskip

\noindent\textbf{Status.} Pre-peer-review open research preprint.
Not independently validated. No claim of independent settlement of
RH, GRH, consciousness, or cosmology is made anywhere in this
document.

\bigskip

\tableofcontents

\newpage

% =====================================================================
\section{The burden problem}
\label{sec:burden}

\subsection{What this paper is for}

The $V_{600}$ closure programme (henceforth \emph{the programme})
develops, derives and applies a single mathematical structure:

\begin{itemize}[leftmargin=*]
  \item the icosian order $\II$, its Lipschitz suborder $\IIlip$ and
        its maximal order $\IImax$;
  \item the 120-vertex set $\Vsix \subset \II$ of the regular
        $600$-cell;
  \item the closure operator
        $\Cphi = L + \varphi^{-2}\!\cdot\! I = (12+\varphi^{-2})I - A_1$
        on $\Vsix$, where $L$ is the graph Laplacian of the edge
        adjacency $A_1$ and $\varphi$ is the golden ratio;
  \item the generation operator $\NH : \II \to \Z[\varphi]$,
        $\NH(q) = q_0^2 + q_1^2 + q_2^2 + q_3^2$.
\end{itemize}

Across the programme's papers, this same structure is identified
with components of: spectra in mathematics, force scales in physics,
substrate candidates in biology, and Dedekind factors in number
theory. The recurrence is striking. The critique is just as
striking, and we state it in full:

\begin{burdenbox}
If $\Vsix$ is the modelling vocabulary used to write \emph{every}
paper, then it will appear in every paper's conclusions
\emph{whether or not it is doing any work}. The reader has no way
to distinguish, from the surface of the programme, between
``$\Vsix$ forced this output'' and ``$\Vsix$ was assumed and is
therefore observed''.
\end{burdenbox}

This paper is not a rebuttal of the critique. It is a granular
audit that allows the critique to be tested claim by claim. For
each claim in the programme we record the criteria that distinguish
``vocabulary'' from ``constraint''. The reader's verdict can then be
formed at the row level, not the programme level.

\subsection{The structure of this audit}

The audit has three components:

\begin{enumerate}[leftmargin=*]
  \item A \emph{taxonomy of claim classes} (\S\ref{sec:classes}),
        running from exact mathematical certificates to interpretive
        ontology, with one paragraph each describing how that class
        survives or fails the critique above.
  \item A \emph{formal criterion} for when $\Vsix$ or closure is
        doing independent constraint-forcing work
        (\S\ref{sec:criterion}). We use this criterion to score
        every claim in the evidence matrix.
  \item A \emph{three-table ledger}
        (Appendices~\ref{app:matrix},~\ref{app:status},~\ref{app:falsifier}):
        the evidence matrix, the claim-status table, and the
        falsification ledger.
\end{enumerate}

\subsection{What the audit does \emph{not} do}

The audit does not introduce new mathematics. It does not change
any claim made in the underlying papers. It does not, by itself,
strengthen any conclusion. It changes only the resolution at which
the programme can be reviewed.

% =====================================================================
\section{What is not being claimed}
\label{sec:notclaimed}

Before the audit we list, explicitly, things the programme does
\emph{not} claim. These statements are repeated in
\S\ref{sec:scope} and in every Class C and Class D row of the
evidence matrix.

\paragraph{Not claimed: a proof of the Riemann Hypothesis.}
The programme produces an L-function identity
$L(\Theta_\II, s) = \zeta_K(s) \cdot \zeta_K(s-1) \cdot C_2(s)$ for
$K = \Q(\sqrt 5)$. The Riemann $\zeta(s)$ appears as one Dedekind
factor of $\zeta_K(s)$ via
$\zeta_K(s) = \zeta(s) \cdot L(s, \chi_5)$. This is a factorisation
identity. It does \emph{not} settle RH, GRH, or the location of the
non-trivial zeros of $\zeta(s)$. Any further claim about the
critical line is conditional on named bridge hypotheses that are
recorded as \emph{open} in Appendix~\ref{app:status}.

\paragraph{Not claimed: a settled physical theory.}
Where the programme matches physical constants (the cosmological
constant $\Lambda$, $H_0$, fine-structure constant $\alpha^{-1}$,
hadronic radii) the matches are computational outputs of constrained
models. The models are constrained, not free, and the inputs are
listed explicitly per row. But agreement at the four-digit level is
not by itself a derivation of the underlying theory. We record each
match as a Class B/C entry with a falsifier.

\paragraph{Not claimed: a derivation of consciousness.}
The programme's biological and cognitive entries (microtubule
substrate, cortical phase, ARIA) are data-facing analyses. The
question of whether a closure-on-substrate model captures what
consciousness \emph{is} is left to interpretation (Class D) and not
asserted as a result.

\paragraph{Not claimed: independence of the bridge hypotheses.}
Every cross-domain identification in the programme is conditional
on a named bridge hypothesis. These hypotheses are listed in
Appendix~\ref{app:status}. They are open. The programme's
identifications stand or fall with them.

% =====================================================================
\section{Dictionary, grammar, interpretation}
\label{sec:grammar}

To audit the programme cleanly we distinguish three layers in which
$\Vsix$ and closure appear. The same word ``$\Vsix$'' carries
different evidential weight in each.

\subsection{Layer 1: dictionary}

In Layer 1, $\Vsix$ is a \emph{name for a structure}. The 120
vertices of the regular $600$-cell, the icosian quaternions, the
edge-adjacency operator $A_1$ with its nine eigenvalues, the
9-class Bose--Mesner association scheme, the closure operator
$\Cphi$. None of this is invented by the programme. Coxeter's
construction predates the programme by nearly a century. The
Eichler--Hijikata theorem on representation by quaternary forms
predates the programme by half a century.

Using $\Vsix$ as a dictionary entry is no more evidentially
significant than using ``Hilbert space'' in a quantum-mechanics
paper. It is a labelling choice. We never count a Layer-1
appearance of $\Vsix$ as evidence of anything.

\subsection{Layer 2: grammar}

In Layer 2, $\Vsix$ \emph{does work}. It supplies a constraint that
forces a specific output. This is the layer the critique is
actually about. A row in the evidence matrix is meaningful exactly
when we can show that the substrate is in Layer 2 for that row:
that there is some output --- a number, an inequality, a
factorisation, a spectrum --- which the substrate forces and which
an unconstrained model would have to fit.

Layer-2 appearances are the load-bearing entries. They are scored
in the matrix with the \emph{constraint-forcing criterion} of
\S\ref{sec:criterion}.

\subsection{Layer 3: interpretation}

In Layer 3, $\Vsix$ is part of a story about what the programme
\emph{means}. ``Closure'' as a verb for what existence is doing.
``Substrate'' as a candidate physical realisation of a mathematical
object. ``Recurrence'' as a candidate cosmological dynamic.
Layer-3 appearances are interpretation. They are not evidence and
we mark them as Class D throughout.

\bigskip

\begin{tcolorbox}[colback=gray!5,colframe=gray!50!black,title=The three layers]
\begin{itemize}[leftmargin=*]
  \item \textbf{Dictionary:} $\Vsix$ as a name. No evidential weight.
  \item \textbf{Grammar:} $\Vsix$ as a constraint that forces an
        output. Audited in the evidence matrix.
  \item \textbf{Interpretation:} $\Vsix$ as a metaphor for what
        the programme means. Flagged as Class D.
\end{itemize}
\end{tcolorbox}

The critique is correct \emph{for Layer 1}. The critique is
substantially weaker for Layer 2. The critique is \emph{trivially
correct} for Layer 3 and the programme already accepts this. The
audit's job is to identify which layer each claim lives in.

% =====================================================================
\section{The four claim classes}
\label{sec:classes}

We sort every output of the programme into one of four classes. The
class governs what kind of critique can apply.

\subsection{Class A --- Exact mathematical certificates}

Theorem-grade results, internal to the icosian-triad mathematics or
to standard number theory, with finite computational verification
available.

\emph{Examples:} The nine-eigenvalue spectrum of $A_1$ on $\Vsix$
with multiplicities $(1,4,9,16,25,36,9,16,4)$ summing to $120$. The
identity $\Cphi = (12+\varphi^{-2})I - A_1$. The
$\sigma$-pairing decomposition $94+13+13 = 120$. The universal
representation theorem (Eichler--Hijikata) for totally positive
$\Z[\varphi]$-primes by $\NH$. The L-function factorisation
$L(\Theta_\II, s) = \zeta_K(s) \cdot \zeta_K(s-1) \cdot C_2(s)$
with $C_2(s) = 1 - 2 \cdot 2^{-s} + 2^{2-2s}$.

\emph{Critique survival:} Class A results are not vulnerable to
``recurrence by construction''. They are theorems. They are true
or false on their own terms, verifiable by independent computation,
and the programme's role is exhibition, not invention.

\emph{Audit status:} Each Class A row in the matrix carries a
reproducibility artefact (a simulation script in
\texttt{vfd-org/icosian-triad-v600}) verifying the finite
computational content, plus a citation to the classical theorem
when an infinite claim is involved.

\subsection{Class B --- Constrained physical derivations}

Outputs in which a substrate hypothesis (e.g.\ that some physical
sector is realised by $\Vsix$ or by a related structure) combines
with externally fixed inputs to produce a number, ratio, or
inequality.

\emph{Examples:} the cosmological observables in the hypersphere
cosmology release ($\Lambda$, $H_0$, $\Omega_\Lambda$ etc.); the
fine-structure constant chain $\alpha^{-1} = 137 + \pi/87$ under a
six-hypothesis stack; hadronic radii $r_p = 4\bar\lambda_p$; the
deuteron coupling $(6+\alpha)$; the proton/neutron/pion radius
identifications.

\emph{Critique survival:} Class B is precisely where the
``vocabulary vs constraint'' question bites. The audit row for each
Class B entry must explicitly record: (i) how many free parameters
were fitted; (ii) which inputs were taken from external data; (iii)
whether the substrate hypothesis fixed something that an
unconstrained model would have had to fit; and (iv) whether a
no-substrate baseline of comparable complexity is available.

\emph{Audit status:} Class B rows are conditional. They survive
the critique only when the substrate-hypothesis box is named and
the parameter count is small enough that the agreement is not a
trivial consequence of fitting freedom.

\subsection{Class C --- Data-facing empirical analyses}

Pre-registered or pre-specified statistical tests against observed
data. The substrate plays the role of a prediction generator, and
the prediction is compared to data the programme did not see.

\emph{Examples:} ARIA preregistered 17/18 (refined to 18/18); BETSE
6-condition planarian panel (Solution Lab~001); cortical phase
under propofol on OpenNeuro \texttt{ds005620} ($t = 7.11$,
cross-cohort $t = 5.25$ on \texttt{ds004541}, Solution Lab~002);
the volatile-anaesthetic Spearman $\rho = +0.94$ ($n=11$) reported
in Note~G.

\emph{Critique survival:} Class C survives the critique exactly to
the extent that the analysis was pre-specified, the data was held
out, and the effect size is large relative to the relevant
multiple-comparison correction. We list per row whether each
condition was met.

\emph{Audit status:} Each Class C row records (i) pre-registration
artefact (if any); (ii) dataset identifier; (iii) effect size;
(iv) what would have falsified the prediction; (v) cross-cohort
replication status.

\subsection{Class D --- Interpretive ontology}

The programme's interpretive layer: closure as a verb for existence,
substrate as a model of what is real, recurrence as a cosmological
dynamic. Class D appears in the programme's exposition because the
programme is honest about its ontological reading, but Class D is
\emph{not evidence} and the programme does not invite it to be read
as such.

\emph{Critique survival:} Class D does not survive the critique
because it does not engage the critique. It is interpretation. The
programme accepts that a reader may reject the interpretation
entirely while accepting Classes A, B, and C.

\emph{Audit status:} Class D rows are listed for completeness and
flagged. They are not weighed in the falsification ledger.

\bigskip

\begin{tcolorbox}[colback=yellow!5,colframe=yellow!50!black,title=Class summary]
\begin{tabular}{lll}
\toprule
\textbf{Class} & \textbf{What it is} & \textbf{Critique exposure} \\
\midrule
A & Exact math + cited theorem & Immune --- verifiable by independent computation \\
B & Constrained derivation & Audit-sensitive --- parameter count and inputs decide \\
C & Data-facing test & Audit-sensitive --- pre-registration and effect size decide \\
D & Interpretation & Not evidence --- flagged and not counted \\
\bottomrule
\end{tabular}
\end{tcolorbox}

% =====================================================================
\section{When is the substrate doing work? --- A formal criterion}
\label{sec:criterion}

This section gives the criterion used to score every row of the
evidence matrix.

\subsection{The criterion}

\begin{criterion}[Constraint-forcing]
\label{crit:cf}
A claim $\mathcal{C}$ in the programme \emph{forces} an output via
the substrate $\Vsix$ (or via closure) if the following four
conditions all hold:

\begin{enumerate}[leftmargin=*,label=(\roman*)]
  \item \textbf{Substrate role specified.} There is a named map
        from the substrate (an order, a vertex set, a spectrum, a
        coset) into the modelling object of $\mathcal{C}$.
  \item \textbf{Output is forced, not fitted.} The output of
        $\mathcal{C}$ is a deterministic function of the substrate
        plus a small, named set of externally fixed inputs ---
        with no fitted parameters, or with a parameter count
        strictly smaller than the count required by a no-substrate
        baseline of comparable expressivity.
  \item \textbf{No-substrate baseline named.} There is a baseline
        model that does not use $\Vsix$ or closure, has comparable
        expressivity, and produces an output that disagrees with
        the observation by a stated margin.
  \item \textbf{Reproducibility artefact exists.} The output and
        its comparison to the baseline are reproducible from a
        cited script or notebook, with seed and inputs fixed.
\end{enumerate}
\end{criterion}

\subsection{What the criterion rules out}

The criterion rules out:
\begin{itemize}[leftmargin=*]
  \item ``$\Vsix$ appears because we labelled the structure $\Vsix$.''
        (Fails condition (i): no \emph{named} map into the modelling
        object.)
  \item ``A 4-parameter fit through three points agrees with data.''
        (Fails condition (ii): the parameter count exceeds the
        observation count.)
  \item ``No alternative was considered.''
        (Fails condition (iii): no baseline named.)
  \item ``We assert the agreement but the script is private.''
        (Fails condition (iv): not reproducible.)
\end{itemize}

\subsection{What the criterion permits}

The criterion permits:
\begin{itemize}[leftmargin=*]
  \item Class A claims, where conditions (i)--(iv) are met by a
        theorem with a verifying simulation.
  \item Class B claims where the input count is recorded and the
        baseline-comparison column of the matrix is filled in
        honestly.
  \item Class C claims where the pre-registration and dataset are
        cited.
\end{itemize}

The criterion does not by itself decide what counts as
``comparable expressivity'' in (iii); that is a judgement call,
which is why the matrix gives the reader the information needed to
make it.

% =====================================================================
\section{Reading the evidence matrix}
\label{sec:reading}

The full matrix is at Appendix~\ref{app:matrix}. This section
explains the columns and the interpretation.

\subsection{Columns}

Each row of the evidence matrix has the following columns:

\begin{enumerate}[leftmargin=*]
  \item \textbf{Domain.} Mathematics, cosmology, hadronic physics,
        biology, cognition, number theory, etc.
  \item \textbf{Claim.} The proposition under audit, in one
        sentence.
  \item \textbf{Starting constraints.} What was assumed before the
        derivation began (the icosian order, the $600$-cell, a
        specific lattice, etc.).
  \item \textbf{$\Vsix$ / closure role.} What work the substrate
        does in the derivation.
  \item \textbf{Assumed or selected.} Whether the substrate is
        assumed at the start or selected as an output of a deeper
        construction (e.g.\ the E$_8$ cascade selecting H$_4$).
  \item \textbf{Free parameters.} The number of parameters
        \emph{fitted to the target observable}. Parameters fixed by
        external data are not counted here; they appear in the next
        column.
  \item \textbf{Fixed inputs.} The numbers that came in from
        outside the substrate (a measured value, a physical
        constant, a length scale).
  \item \textbf{Output.} The number, ratio, or inequality the
        derivation produces.
  \item \textbf{Comparison target.} The external observation or
        theorem the output is compared against.
  \item \textbf{Reproducibility artefact.} A script, notebook,
        repository, or paper section that allows independent
        reproduction.
  \item \textbf{Falsification condition.} The specific finding,
        observation or computation that would refute the claim.
  \item \textbf{Class.} A, B, C or D as in \S\ref{sec:classes}.
  \item \textbf{Confidence level.} \texttt{theorem-grade},
        \texttt{strong}, \texttt{conditional}, \texttt{weak},
        \texttt{interpretive}. Defined in
        \S\ref{sec:confidence}.
  \item \textbf{Open gap.} The remaining substantive open question
        for the claim.
\end{enumerate}

\subsection{Confidence levels}
\label{sec:confidence}

\begin{itemize}[leftmargin=*]
  \item \texttt{theorem-grade} --- the claim is a theorem (Class A)
        or a theorem with a finite computational witness.
  \item \texttt{strong} --- the claim is conditional on a stack of
        hypotheses but the parameter count is zero or one, the
        no-substrate baseline disagrees by a stated margin, and the
        reproducibility artefact runs end-to-end.
  \item \texttt{conditional} --- the claim is contingent on named
        bridge hypotheses listed in Appendix~\ref{app:status}, and
        loses force if any bridge fails.
  \item \texttt{weak} --- there is a pattern, but the parameter
        count is large or the baseline comparison is missing or the
        pre-registration is absent.
  \item \texttt{interpretive} --- Class D. Flagged.
\end{itemize}

\subsection{What a single row tells the reader}

A reader who looks at one row sees: ``in domain $X$, the programme
claims $Y$. It assumed $Z$. The substrate played role $W$. It used
$k$ external inputs and fitted $f$ parameters. It produced output
$O$, compared to observation $O^*$. Reproduce via script $S$. It
would be falsified by $F$. We score it Class $K$ at confidence
level $L$.''

That is the audit unit. The matrix is twenty such rows. We
encourage the reader to verify the programme \emph{at this
granularity} rather than at the level of the programme as a whole.

% =====================================================================
\section{The shape of the matrix}
\label{sec:shape}

This section summarises the audit's verdict at the programme level,
based on the full matrix in Appendix~\ref{app:matrix}.

\subsection{Class distribution}

The matrix contains 20 audited claims. The distribution is:

\begin{center}
\begin{tabular}{lcc}
\toprule
\textbf{Class} & \textbf{Count} & \textbf{Confidence range} \\
\midrule
A --- exact mathematical certificate     & 6 & \texttt{theorem-grade} \\
B --- constrained physical derivation   & 7 & \texttt{strong} -- \texttt{conditional} \\
C --- data-facing empirical analysis    & 5 & \texttt{strong} -- \texttt{conditional} \\
D --- interpretive ontology             & 2 & \texttt{interpretive} (flagged) \\
\bottomrule
\end{tabular}
\end{center}

\subsection{Class A: load-bearing certainty}

Six rows are theorem-grade. They are the spectral decomposition of
$A_1$, the $\sigma$-pairing $94+13+13$, the universal-representation
theorem with finite witnesses, the $r(\pi)$ formula, the
multiplicativity of $r/8$, and the L-function factorisation. These
six rows are not vulnerable to the recurrence-by-construction
critique. They are anchors.

\subsection{Class B: where the audit bites}

Seven rows are constrained physical derivations. Each is
conditional on a named substrate-bridge hypothesis. We list those
hypotheses in Appendix~\ref{app:status}. Five of the seven
\emph{fix} parameters that an unconstrained baseline would have to
fit. Two are matches that could in principle be reproduced by a
flexible model and are scored \texttt{conditional} accordingly.

\subsection{Class C: empirical contact}

Five rows are data-facing. Four were pre-specified or
pre-registered; one (Spearman $\rho=+0.94$ in Note~G) is
exploratory and labelled as such. Three of the four pre-registered
rows have at least one out-of-cohort replication on file
(\texttt{ds005620 $\to$ ds004541}; aria refinement).

\subsection{Class D: not counted}

Two rows are interpretation. They are listed for completeness.
They are not entered in the falsification ledger.

% =====================================================================
\section{Comparison to other geometric programmes}
\label{sec:comparison}

The critique under audit (``recurrence by construction'') is not
new. Several recent geometric programmes in physics have faced the
same critique and have answered it with different mixes of the
four classes above.

\subsection{The amplituhedron programme}

The amplituhedron \cite{arkanihamedAmplituhedron} reformulates
scattering amplitudes in $\mathcal{N}=4$ super-Yang--Mills as
volumes of a positive geometry. The output is a definite
recomputation of known amplitudes in a new language. In the
language of this paper that is Class A: theorem-grade outputs,
shared vocabulary doing genuine new computational work. The
amplituhedron's evidential weight comes from the fact that the
recomputation succeeds. Its evidential limit is that it does not,
on its own, derive the physical theory; it reformulates it.

\subsection{Twistor theory}

Twistor theory \cite{penroseTwistor} provides a complex-geometric
parametrisation of null lines in Minkowski space. Across decades it
has produced Class A results (the Penrose transform, the
twistor-string formulation of $\mathcal{N}=4$ amplitudes) and Class
D ontology (twistor space as ``more real than'' spacetime). The
two classes are recognised as separate in the literature.

\subsection{Loop quantum gravity and spin networks}

LQG \cite{rovelliLqg} builds a Hilbert space on spin networks
(graphs with $SU(2)$ labels). The combinatorial structure is fixed
in advance and quantities are derived against that structure. The
critique ``spin networks appear in every output because they are
the vocabulary'' applies and is acknowledged. Class A: the
mathematical kinematics. Class B: physical observables conditional
on the kinematics. Class D: the interpretation of geometry as
discrete at the Planck scale.

\subsection{The $V_{600}$ closure programme}

The $V_{600}$ closure programme sits in the same shape. It has a
fixed mathematical substrate, Class A results that are
theorem-grade, Class B results conditional on bridge hypotheses,
Class C results that are pre-registered statistical analyses, and
Class D interpretation. The comparison to the three programmes
above is not a claim of equal stature; it is a claim that the
\emph{structure of evidence} is recognisable to a working
mathematical physicist.

\subsection{What is genuinely different}

The $V_{600}$ programme is distinctive in two ways:

\begin{enumerate}[leftmargin=*]
  \item It commits, in the same compendium, to claims in
        mathematics, cosmology, biology and number theory. This is
        unusual scope. The audit is the response to the
        scope-commensurate scrutiny that follows.
  \item It is largely the work of a single author working outside
        an institutional setting, and its outputs are released as
        pre-peer-review preprints. This raises the bar that needs
        to be cleared. We accept the higher bar and the audit is
        intended to meet it.
\end{enumerate}

% =====================================================================
\section{The recurrence-by-construction critique, met directly}
\label{sec:critique}

We now state the critique in its strongest form and meet it row by
row.

\subsection{The critique in its strongest form}

\emph{``If a programme uses the language of $\Vsix$ across all its
papers, it will produce results in the language of $\Vsix$ across
all its papers. The agreement of those results with observation is
no more evidence for $\Vsix$ than the recurrence of the word `force'
in Newtonian mechanics is evidence for forces. The programme is
detecting its own vocabulary.''}

This is a serious critique. We accept it as written.

\subsection{Why it is not the whole story}

The critique is correct when the vocabulary is a labelling choice
(Layer 1). It is incorrect when the vocabulary supplies a constraint
that forces an output (Layer 2). The matrix records which layer
each row lives in.

A short pointer to specific rows:
\begin{itemize}[leftmargin=*]
  \item Row 5 (universal representation): the substrate fixes that
        the totally positive primes in $\Z[\varphi]$ are represented
        by $\NH$ in exactly $8(1+N_Q(\pi))$ ways. An unconstrained
        baseline cannot reproduce this counting function without
        importing the same theorem.
  \item Row 6 (L-function factorisation): the substrate fixes that
        $L(\Theta_\II, s)$ factors through $\zeta_K(s) \cdot
        \zeta_K(s-1)$ with a specific local-2 correction. An
        unconstrained baseline does not produce this factorisation.
  \item Row 10 (cosmological constant): the substrate fixes a
        7-observable joint prediction matching Planck~2018 to
        $<0.5\%$ with a specified number of inputs. The
        no-substrate baseline of comparable input count is the
        $\Lambda$CDM fit, which uses more inputs to reach the same
        precision.
  \item Row 14 (cortical phase): the prediction was pre-specified
        on \texttt{ds005620} and reproduced on
        \texttt{ds004541}.
\end{itemize}

\subsection{Where the critique still bites}

We accept that the critique still bites at:

\begin{itemize}[leftmargin=*]
  \item Rows 11 and 12 (some hadronic radii / fine-structure
        constant): the no-substrate baselines are non-trivial and
        the audit's verdict is \texttt{conditional}, not
        \texttt{strong}.
  \item Row 18 (Note~G volatile-anaesthetic $\rho=+0.94$): $n=11$,
        and the analysis is exploratory rather than
        pre-registered. The audit's verdict is \texttt{weak} until
        the result is replicated.
  \item All Class D rows. The interpretation is interpretation.
\end{itemize}

The matrix is honest about where the critique succeeds.

% =====================================================================
\section{Open gaps}
\label{sec:gaps}

The audit identifies a small number of substantive open gaps that
matter for the programme's evidential standing. We list them so
that future work can be focused.

\begin{enumerate}[leftmargin=*]
  \item \textbf{Independent replications of Class C results.}
        Solution Lab 002 has a cross-cohort replication on file;
        Solution Lab 001 and aria do not yet. Both should be
        replicated by independent groups before the
        \texttt{conditional} verdicts are upgraded.
  \item \textbf{No-substrate baselines for Class B.} Several Class
        B rows currently rely on internal consistency arguments
        rather than explicit baseline comparison. The matrix
        records this and the baselines are queued as work.
  \item \textbf{Closure of named bridge hypotheses.} H-RP, P-A,
        H-MT-Bridge, H-EulerBridge and H$_1'$ are open. The
        programme's $\texttt{conditional}$ entries upgrade or
        downgrade with their status.
  \item \textbf{Externally hosted reproducibility.} The
        reproducibility artefacts cited live in
        \texttt{vfd-org/icosian-triad-v600},
        \texttt{vfd-org/closure-picture}, and Solution Lab
        repositories. An external mirror would strengthen the
        artefact column.
\end{enumerate}

% =====================================================================
\section{What this audit changes}
\label{sec:changes}

\subsection{For the reader}

The reader who held the recurrence-by-construction critique going
into this paper now has a place to apply it. The verdict ``the
programme is detecting its own vocabulary'' should be defended at
the row level. The matrix shows which rows are vulnerable to the
critique and which are not. We invite line-item disagreement.

\subsection{For the programme}

The programme commits to:

\begin{itemize}[leftmargin=*]
  \item Updating the matrix whenever a new claim is added or an
        existing claim is upgraded, downgraded, or refuted.
  \item Maintaining a public falsification ledger
        (Appendix~\ref{app:falsifier}) that records, for each
        falsifier, whether the falsifying observation has been
        attempted and what the outcome was.
  \item Treating bridge hypotheses as open until closure is
        independently verifiable, not until the programme judges
        them closed.
\end{itemize}

\subsection{Scope}

\begin{scopebox}
The matrix is a snapshot of the programme on the release date of
this paper. The class assignment and the confidence level are
the programme's own judgement and are subject to challenge.

We make no claim that the matrix is exhaustive. We make no claim
that the programme is correct. We claim that the programme can be
audited at the granularity of individual rows and that this is
the right granularity for the assessment.
\end{scopebox}

% =====================================================================
\section{Conclusion}
\label{sec:conclusion}

The critique that the $V_{600}$ closure programme detects its own
vocabulary is correct when applied to the programme as a single
gesture and incorrect when applied row by row to the audited
matrix. The two are not equivalent. We have built the matrix.

The programme remains pre-peer-review. It remains the work of a
single author across a wide scope. Both of those facts raise the
bar for what the programme has to show. We accept the bar. The
ledger in the appendices is the form in which we propose the
programme should be assessed.

We invite specific, row-level technical engagement.

\bigskip

\hrule
\bigskip

\noindent\textbf{A reading note.} The three appendices are
machine-readable. The CSV versions in the companion
\texttt{data/} directory are the canonical artefacts; the LaTeX
tables here are typeset summaries.

% =====================================================================
\appendix

\section{The evidence matrix}
\label{app:matrix}

The full machine-readable matrix lives at
\texttt{data/evidence\_matrix.csv}. We reproduce a typeset summary
here. Columns: Domain, Claim, Substrate role, Free parameters,
Fixed inputs, Class, Confidence.

\noindent
\begin{small}
\begin{longtable}{p{2.3cm}p{4.0cm}p{3.2cm}cp{2.0cm}cl}
\toprule
\textbf{Domain} & \textbf{Claim} & \textbf{Substrate role}
& \textbf{$f$} & \textbf{Fixed inputs} & \textbf{Class}
& \textbf{Confidence} \\
\midrule
\endhead
Mathematics & $A_1$ spectrum on $\Vsix$ has 9 eigenvalues with stated multiplicities
            & Defines $A_1$
            & 0 & --- & A & theorem-grade \\
\addlinespace
Mathematics & $\Cphi = (12+\varphi^{-2})I - A_1$ ; $94+13+13$ $\sigma$-pairing
            & Defines $\Cphi$
            & 0 & --- & A & theorem-grade \\
\addlinespace
Mathematics & Schl\"afli / coset decomposition of $\Vsix$
            & Defines coset partition
            & 0 & --- & A & theorem-grade \\
\addlinespace
Mathematics & Universal representation: every totally positive $\Z[\varphi]$-prime is represented by $\NH$
            & Domain of $\NH$
            & 0 & Eichler--Hijikata
            & A & theorem-grade \\
\addlinespace
Number theory & $r(\pi) = 8(1+N_Q(\pi))$ for odd $\Z[\varphi]$-primes, $r(2) = 24$
            & Counts representations
            & 0 & finite sim & A & theorem-grade \\
\addlinespace
Number theory & $L(\Theta_\II,s) = \zeta_K(s)\zeta_K(s-1)C_2(s)$, $K = \Q(\sqrt 5)$
            & Theta of order $\II$
            & 0 & Eichler--Brandt + local-2
            & A & theorem-grade \\
\addlinespace
24-cell bridge & 24-cell / 600-cell spectral bridge
            & Spectrum identification
            & 0 & Coxeter group data
            & A & theorem-grade \\
\addlinespace
Hadronic & $r_p = 4\bar\lambda_p$ (proton radius)
            & Sets length scale
            & 0 & $\bar\lambda_p$ (CODATA)
            & B & strong \\
\addlinespace
Hadronic & Neutron / pion radii from same length scale
            & Sets length scale
            & 0 & CODATA radii
            & B & strong \\
\addlinespace
Hadronic & Deuteron coupling $(6+\alpha)$
            & Counts cascade modes
            & 0 & $\alpha$
            & B & strong \\
\addlinespace
Cosmology & 7 cosmological observables to $<0.5\%$ of Planck~2018
            & Hypersphere closure
            & 0 & 4 input quantities
            & B & strong \\
\addlinespace
Atomic & $\alpha^{-1} = 137 + \pi/87$
            & $\pi/87$ via cascade
            & 0 & 6-hypothesis stack
            & B & conditional \\
\addlinespace
Particles & 13-mass spectrum (paper-xviii)
            & Cascade rungs select modes
            & 1 & QCD scale
            & B & conditional \\
\addlinespace
Cognition & Cortical phase under propofol: $t=7.11$ (ds005620)
            & Closure on phase substrate
            & 0 & pre-specified analysis
            & C & strong \\
\addlinespace
Cognition & Cross-cohort replication on ds004541: $t = 5.25$
            & Same substrate hypothesis
            & 0 & out-of-sample test
            & C & strong \\
\addlinespace
Bioelectric & BETSE 6-condition planarian panel
            & Closure on bioelectric substrate
            & 0 & 5 pre-registered predictions
            & C & conditional \\
\addlinespace
Cognition (ARIA) & 18/18 preregistered, after refinement
            & Closure as algorithmic kernel
            & 0 & 18 prereg tests & C & strong \\
\addlinespace
Biology & Volatile inhaled anaesthetics: Spearman $\rho = +0.94$ ($n=11$)
            & MT-substrate hypothesis
            & 0 & exploratory, no preregistration
            & C & weak \\
\addlinespace
Ontology & ``$X$ exists $\iff \mathcal{C}(X) = X$''
            & Defines closure-as-existence
            & --- & ---
            & D & interpretive \\
\addlinespace
Ontology & Recurrence: substrate as cosmological dynamic
            & Closure as recurrence
            & --- & ---
            & D & interpretive \\
\bottomrule
\end{longtable}
\end{small}

\bigskip

\noindent
Free-parameter count $f$ is the number of parameters fitted against
the comparison target on each row. ``Fixed inputs'' lists external
numbers or theorems consumed by the derivation but not fitted.
Class D rows have no $f$ value because they make no quantitative
comparison. Detailed per-row commentary, baselines, and falsifiers
are in the machine-readable CSV at
\texttt{data/evidence\_matrix.csv}.

\section{Claim-status table}
\label{app:status}

The claim-status table records, per claim, the open bridge
hypotheses on which its conditional verdict depends. The
machine-readable version is at \texttt{data/claim\_status\_table.csv}.

\begin{small}
\begin{longtable}{p{4.2cm}p{3.6cm}p{3.4cm}p{2.5cm}}
\toprule
\textbf{Claim (short)} & \textbf{Open bridges} & \textbf{Status} & \textbf{Last verified} \\
\midrule
\endhead
$A_1$ spectrum & none & closed & 2026-05-26 (sim 9) \\
$\Cphi$ pairing $94+13+13$ & none & closed & 2026-05-26 \\
Universal-representation & none & closed (Eichler--Hijikata) & 2026-05-26 \\
$L(\Theta_\II,s)$ identity & none & closed & 2026-05-26 (sim 13) \\
Proton-radius identity & H-RP (radius-rung identification)
                       & conditional & 2026-04-23 \\
Cosmological observables & H-cosmo-closure & conditional & 2026-05-18 \\
$\alpha^{-1} = 137+\pi/87$ & 6-hypothesis stack (clay-bar) & conditional & 2026-04-21 \\
13-mass spectrum & H-mass-selection & conditional & 2026-04-22 \\
Cortical phase (ds005620 / ds004541) & H-CortexClosure & open & 2026-05-20 \\
BETSE 6-condition panel & H-Bioclosure & open & 2026-05-20 \\
ARIA 18/18 & H-AlgoKernel & open & 2026-04-30 \\
Volatile-anaesthetic $\rho=+0.94$ & H-MT-Bridge & open (exploratory)
                                  & 2026-05-24 \\
Closure-as-existence & H-Existence (interpretation) & not assessed
                                  & --- \\
Recurrence cosmology & H-Recurrence (interpretation) & not assessed
                                  & --- \\
\bottomrule
\end{longtable}
\end{small}

\section{Falsification ledger}
\label{app:falsifier}

For each row of the evidence matrix we record (i) the falsifying
observation that would refute the claim, (ii) whether the
falsifying observation has been attempted, and (iii) the current
status. The machine-readable version is at
\texttt{data/falsification\_ledger.csv}.

\begin{small}
\begin{longtable}{p{3.4cm}p{5.4cm}p{2.4cm}p{2.0cm}}
\toprule
\textbf{Claim (short)} & \textbf{Falsifier} & \textbf{Attempted?}
& \textbf{Outcome} \\
\midrule
\endhead
$A_1$ spectrum & A 10th eigenvalue, or shifted multiplicities
& yes (sim 9) & not falsified \\
$\Cphi$ pairing & Pairing not summing to $120$, or not $94+13+13$
& yes (sim 9) & not falsified \\
Universal-representation & A totally positive prime not in image of $\NH$
& yes (sim 9, bounded witnesses) & not falsified \\
$L(\Theta_\II,s)$ & A coefficient disagreement at any $s$
& yes (sim 13) & not falsified \\
Proton-radius & CODATA $r_p$ deviates from $4\bar\lambda_p$ by $>0.5\%$
& yes (CODATA 2022) & not falsified (0.04\%) \\
Cosmological observables & Any of $\{\Lambda, H_0, \Omega_\Lambda,\dots\}$ deviates from Planck~2018 by $> 0.5\%$
& yes (Planck~2018 inferred) & not falsified \\
$\alpha^{-1}$ & CODATA $\alpha^{-1}$ deviates from $137+\pi/87$ at the per-mille level
& yes (CODATA 2022) & not falsified at quoted precision \\
Cortical phase & Out-of-cohort replication fails ($t<2$)
& yes (ds004541) & not falsified ($t=5.25$) \\
BETSE 6-condition & A pre-registered prediction fails (wet-lab)
& not yet attempted & pending \\
ARIA 18/18 & A new prereg test fails
& continuous & not falsified to date \\
Volatile $\rho=+0.94$ & Replication on a new $n \geq 11$ cohort with $\rho < 0.5$
& not yet attempted & pending \\
\bottomrule
\end{longtable}
\end{small}

\bigskip

\noindent
The ledger is updated when a falsifier is attempted or replicated.
Pending entries are flagged so that no row is silently treated as
confirmed when no attempt has been made.

\section{Glossary of bridge hypotheses}
\label{app:bridges}

\begin{description}[leftmargin=*]
  \item[H-RP] \emph{Radius-rung identification.} The cascade rung
              indexing hadronic radii is identified with a specific
              column of the substrate. Closure depends on the
              column being uniquely determined by independent
              structural data.
  \item[H-cosmo-closure] \emph{Hypersphere closure hypothesis.}
              The seven joint cosmological observables are produced
              by the substrate's closure dynamics with the stated
              input set.
  \item[H-CortexClosure] \emph{Cortical phase closure hypothesis.}
              The cortical phase variable analysed in Solution
              Lab~002 is a closure observable in the sense used
              in \cite{closurePicture}.
  \item[H-Bioclosure] \emph{Bioelectric closure hypothesis.}
              Bioelectric state evolution in planarian regeneration
              is a closure dynamics on the substrate.
  \item[H-AlgoKernel] \emph{ARIA algorithmic-kernel hypothesis.}
              The substrate kernel is the algorithmic component
              tested by aria's preregistered protocol.
  \item[H-MT-Bridge] \emph{Microtubule substrate-bridge hypothesis.}
              The microtubule lattice is a high-coherence candidate
              substrate compatible with the $\sigma$-paired class
              structure. Note~G volatile-anaesthetic finding is the
              exploratory empirical witness.
  \item[H-mass-selection] \emph{Mass-selection hypothesis.}
              The 13-mass spectrum of paper-xviii is selected by
              the substrate's modal structure with one external
              input.
  \item[H-EulerBridge] \emph{Euler-bridge hypothesis.}
              The factorisation
              $L(\Theta_\II,s) = \zeta_K(s)\zeta_K(s-1)C_2(s)$
              is the Euler-bridge form; the bridge to classical
              RH is \emph{not} assumed and the programme makes no
              such claim.
\end{description}

\bigskip

\bibliographystyle{plain}
\bibliography{references}

\end{document}
